% Manuscript insertion: experimental scope and claim boundaries
We distinguish two evaluation surfaces. The \textbf{matched manuscript surface} is the primary contract used for main-text claims: fixed action budgets, predefined method set, and standardized scoring/evaluation slices. A broader \textbf{operational surface} (held-out checks, dry-run scaffolds, and cross-provider real-model audits) is used for boundary testing and robustness diagnostics.

To avoid selective framing, we report the matched-surface results in the main text and place held-out and real-model analyses in appendix-oriented evidence layers. These auxiliary layers inform claim boundaries (e.g., instability and non-dominance risks), but are not used to assert universal superiority. In particular, this paper does \emph{not} claim universal SOTA performance across providers, datasets, or deployment conditions; instead, it claims an operationally specified adaptive-allocation formulation with bounded empirical support under the stated action-budget contract.

For anti-collapse, claim boundaries are also calibration-bound: we report non-monotonic ablation behavior and the matched-surface calibration sweep as evidence of tradeoff sensitivity, not as universal proof that the default anti-collapse setting should always be kept or always removed.
